\section{Introduction}
\label{sec:introduction}

Roman Urdu---the phonological transcription of the Urdu language using the standard 26-letter Latin alphabet---has emerged as the de facto written medium for informal digital communication across Pakistan and the broader South Asian diaspora. It permeates user reviews on e-commerce platforms, customer feedback on social media channels, and mobile instant messaging. Despite its ubiquity, Roman Urdu is strictly an unstandardized vernacular. While classical Nastaliq Urdu employs an alphabet of 36 distinct characters with formal orthographic rules, Roman Urdu lacks any centralized standardizing authority or agreed-upon phonetic mapping. Consequently, native speakers write purely by acoustic approximation. 

A single lexical item or sentiment-bearing predicate can easily manifest in several divergent spelling variations depending entirely on individual typists, regional phonology, and texting conventions. For instance, the core affirmation \textit{accha} (``good'') frequently appears as \textit{acha}, \textit{achha}, \textit{achaah}, or \textit{achaaa}; similarly, the negation \textit{nahi} (``not'') is written interchangeably as \textit{nhn}, \textit{nahin}, \textit{nhi}, or \textit{nahe}. In the context of automated Natural Language Processing (NLP), and Sentiment Analysis (SA) in particular, this extreme orthographic entropy poses a formidable challenge.

In contemporary NLP research, Pre-trained Multilingual Language Models (PMLMs) such as multilingual BERT (\mbox{mBERT})~\cite{devlin2019bert} and XLM-RoBERTa (\mbox{XLM-R})~\cite{conneau2020unsupervised} have demonstrated state-of-the-art classification performance on code-mixed and low-resource benchmarks. However, the standard evaluation paradigm in the literature remains fundamentally flawed: models are trained and tested on single, static, identically distributed (i.i.d.) random splits of a scraped corpus. When an automated classifier achieves high benchmark accuracy under such conditions, it often reflects memorization of canonical spelling forms within a closed vocabulary rather than generalized sentiment comprehension. In live production environments, incoming text exhibits uncontrolled spelling shifts that induce severe out-of-vocabulary (OOV) rates, token fragmentation, and downstream classification failure.

Recent literature explicitly acknowledges this bottleneck. Notably, Nazir et al.~\cite{nazir2025leveraging} benchmarked mBERT on code-mixed Roman Urdu and Roman Punjabi social media reviews, demonstrating strong in-domain accuracy (79.94\%) but observing that generic text preprocessing (e.g., stop-word and punctuation removal) produced only marginal improvements (2.71\%). The authors highlighted that imperfect Romanized text normalization remains an unaddressed limitation. In parallel, Khan et al.~\cite{khan2026romeva} investigated vocabulary expansion on the Roman Urdu Word Variations and Normalized Sentiment Review (RUWV-NSR) dataset~\cite{ahmed2024ruwvnsr}, focusing on embedding-space drift. However, no prior work has systematically quantified the degradation curves of multilingual transformers under controlled, graded orthographic noise for Roman Urdu, nor investigated whether adaptation strategy influences robustness.

This study directly bridges this empirical gap. Treating Roman Urdu sentiment classification as a rigorous, end-to-end data-science pipeline, we investigate the following core research questions:
\begin{enumerate}
    \item[\textbf{RQ1:}] \textit{Degradation Under Orthographic Noise:} How does sentiment classification performance degrade across increasing severity levels of authentic spelling variation for classical versus transformer-based architectures?
    \item[\textbf{RQ2:}] \textit{Parameter-Efficient vs. Full Adaptation:} Does Low-Rank Adaptation (LoRA)~\cite{hu2022lora} confer greater robustness against orthographic perturbation compared to unconstrained full fine-tuning?
    \item[\textbf{RQ3:}] \textit{Token Fragmentation Profiling:} To what extent do distinct subword tokenization schemes (WordPiece vs. SentencePiece/BPE) fragment authentic lexical variants, and how does this correlate with performance drop?
    \item[\textbf{RQ4:}] \textit{Factorial Disentanglement:} Is performance degradation predominantly driven by word substitution rate, character edit distance, or their mutual interaction?
    \item[\textbf{RQ5:}] \textit{Mitigation and Normalization:} Can lexicon-driven orthographic normalization or noise-augmented training reverse perturbation losses where generic preprocessing fails?
\end{enumerate}

To answer these questions, we perform an independent data audit on the RUWV-NSR corpus, resolving pervasive label corruption and eliminating 4,632 exact duplicate reviews (16.49\%) to establish a strictly leakage-free benchmark. Leveraging a paired lexicon of 5,244 words (10,869 authentic variant pairs), we develop a perturbation engine that constructs canonical noise tiers alongside a full $3 \times 3$ factorial grid. We benchmark four primary modeling paradigms: TF-IDF + Logistic Regression, fully fine-tuned mBERT, fully fine-tuned XLM-RoBERTa, and LoRA-adapted XLM-RoBERTa. Finally, we execute normalization ablations and validate cross-domain generalizability on 3,000 independent reviews from the Daraz.pk RUECD corpus~\cite{chandio2022sentiment}.
